\documentclass[11pt,a4paper]{article}
\usepackage[margin=25mm]{geometry}
\usepackage[T1]{fontenc}
\usepackage{lmodern}
\usepackage{amsmath,amssymb,bm,booktabs,tabularx}
\usepackage[numbers,sort&compress]{natbib}
\usepackage[colorlinks=true,allcolors=blue]{hyperref}
\newcommand{\MSE}{\operatorname{MSE}}
\newcommand{\Peak}{\operatorname{MaxSqErr}}
\newcommand{\todo}[1]{\par\noindent\textbf{To complete:} #1\par}
\title{Learning Compact B-Spline Representations through\\
Candidate Selection and Subset-Conditioned Geometry}
\author{Author names and affiliations to be supplied}
\date{Working manuscript framework --- 21 September 2026}

\begin{document}
\maketitle
\begin{abstract}
Selecting a compact knot vector while preserving approximation accuracy is a
coupled discrete--continuous problem. We describe a learned pipeline that maps
ordered samples to candidate knots, selects a variable-size subset, and jointly
updates retained knot positions and point parameters. A curve-adaptive threshold
controls probability-mass Top-$K$ selection. A subset-conditioned decoder
preserves the correspondence between parameter updates and knot relocation,
and standard B-spline least squares determines the control points.
During training, bounded online counterfactual search supplies verified
feasible subsets and geometric targets; this search is absent at deployment.
We specify evaluation on synthetic splines, handwriting strokes, geographic
curves, and procedural industrial offsets, reporting failures, cardinality,
pointwise peak error, and complete latency. Current diagnostic runs show that
improved fit feasibility does not automatically produce more compact
representations. Established components are therefore separated from planned
depth, capacity, and peak-error ablations, without claiming universal
feasibility or globally minimal knot count.
\end{abstract}
\noindent\textbf{Keywords:} B-spline fitting; knot selection; parameterization;
geometric learning; counterfactual supervision.

\section{Introduction}
Ordered point sequences arise in digitized curves, map polylines, and
manufacturing-related geometry. Dense spline representations reduce residual
error but introduce unnecessary degrees of freedom. Conversely, deleting knots
can invalidate an accurate fit. Selection must therefore be considered together
with parameterization and the positions of retained knots.

Our aim is to amortize part of numerical knot search into a predictor while
retaining an explicit standard B-spline output. The implemented design combines:
\begin{enumerate}
\item local geometric candidate generation and curve-adaptive one-shot selection;
\item survivor-conditioned parameter and knot updates anchored to proposal geometry;
\item online feasible-subset supervision verified through the actual decoder.
\end{enumerate}
These are design components, not established novelty or superiority claims.
The research question is whether learned selection can approach the
compactness of numerical search while retaining accuracy on unseen curves.

\section{Related work}
Park and Lee~\cite{park2007} use dominant points for adaptive refinement.
Liang et al.~\cite{liang2017} investigate initial knot placement and iterative
insertion for error-bounded fitting. Dung and Tjahjowidodo~\cite{dung2017}
combine coarse subdivision with local position and continuity optimization.
Kang et al.~\cite{kang2015} select active knots by sparse optimization and
subsequently refine the representation. Luo et al.~\cite{luo2022} combine an
$\ell_{\infty,1}$ sparse stage and differential evolution. The last reference
is the sparse/DE paper, not another deep-learning paper of the same year.

The repository implements disclosed adaptations, sharing a final-refit and
evaluation protocol but not every native objective, continuity model, or
solver. Comparing these implementations is not an exact reproduction of the
authors' published results.
\todo{Extend the learning-related discussion using verified primary references.
Distinguish amortized search from guaranteed constrained optimization.}

\section{Problem and error definitions}
Let $\bm Q=(\bm q_i)_{i=1}^M$ be normalized ordered observations in
$\mathbb R^d$, with $0=t_1<\cdots<t_M=1$. A clamped cubic spline is
\begin{equation}
\bm C(t)=\sum_{a=1}^{K+4}B_{a,3}(t;\bm U)\bm P_a,\qquad
\bm U=(0,0,0,0,u_1,\ldots,u_K,1,1,1,1).
\end{equation}
$K$ counts interior knots only: $K=16$ means 24 knot-vector entries and
20 control points, while $K=32$ means 40 entries and 36 control points.
Endpoint-constrained least squares computes
$\bm P^\star$ for given $\bm t$ and $\bm U$. Define
\begin{equation}
e_i=\|\bm C(t_i)-\bm q_i\|_2^2,\qquad
\MSE=M^{-1}\sum_i e_i,\qquad \Peak=\max_i e_i.
\end{equation}
Neither metric is square-rooted or divided by coordinate dimension. The goal is
\begin{equation}
\min_{\bm t,\bm U,K}K
\quad\text{subject to}\quad
\MSE(\bm t,\bm U,\bm P^\star)\leq\varepsilon,\quad 0\leq K\leq K_c,
\label{eq:target}
\end{equation}
with ordered parameters and admissible knots. A one-shot prediction approximates
this goal, without certifying feasibility or global minimality. A peak target
$\delta$, if enabled, is separate from $\varepsilon$. Discrete corresponding-point
errors are not Hausdorff distances or guarantees between observations.

\section{One-shot prediction framework}
\subsection{Point features and proposal parameters}
The geometry encoder consumes coordinates and sequential features, including
first and second differences. Pointwise and pooled features feed a ParameterHead,
which predicts monotone parameters relative to chord-length parameters.
Positive interval constraints, endpoint normalization, and a learned trust
factor restrict the update. Parameter positional encodings augment point
features used as attention memory. This head reads encoded observations,
not ground-truth knots.

\subsection{Candidate generation}
Candidate queries use local Gaussian-biased cross-attention to read
parameter-positioned point features. The CandidateKnotHead produces
$K_c$ ordered interior coordinates $\bm U^{(p)}$ and tokens $\bm h_j$.
These are learned proposals, not necessarily exact shape-preserving Boehm
insertions. An explicit redundant input spline is not required at deployment.

\subsection{Adaptive selection and variable cardinality}
Candidate interaction combines point memory, neighborhood coverage features,
and a tolerance embedding. With importance scores $a_j$ and a curve-level
adaptive threshold $\beta$, the profile uses
\begin{equation}
p_j=\sigma(a_j-\bar a-\beta),\qquad m=\sum_{j=1}^{K_c}p_j.
\end{equation}
For zero safety reserve, requested cardinality is zero if $m<0.5$, and
$\lceil m\rceil$ otherwise, clipped to its admissible range. In the current
profile this range is $[4,K_c]$, so a zero raw request still retains at least
four knots. The highest-probability candidate in each of four nonempty
parameter bins is given priority when the requested count can accommodate
all bin anchors; remaining slots follow the probability ranking. KeepMask
is obtained by one stable Top-$K$ operation on this adjusted ranking. Thus independent
$p_j\geq0.5$ tests are \emph{not} this profile's selection rule. Optional
reserve settings and changes to the four-bin coverage rule must be disclosed
separately. The mask supplies
the output length; separate decoders for every possible $K$ are not evaluated.

\subsection{KeepMask-conditioned geometry}
Distances to nearest retained neighbors, retained rank, and retained count
form a geometry embedding added to each candidate token. Here a
\emph{keep embedding} means mask-conditioned geometry, not an independent
binary lookup table. Only retained tokens supply survivor-attention keys and
values. Point queries also attend to this survivor memory to update parameter
features. A masked zero sentinel handles empty subsets. Discarded candidates
are not used as survivor evidence, and candidates are not independent.

\subsection{Anchored parameter--knot coupling}
Let $\widehat{\bm t}$ denote a monotone subset-conditioned parameter update.
The decoder limits its difference from proposal parameters and contracts the
update when required to preserve selected knot gaps:
\begin{equation}
\bm t'=\bm t^{(p)}+\gamma s(\widehat{\bm t}-\bm t^{(p)}),
\qquad 0\leq\gamma,s\leq1.
\end{equation}
The same piecewise-linear parameter warp transforms candidate knots.
A retained warped knot $\widetilde u_j$ is then relocated within free room
$r_j$ to its nearest retained neighbors:
\begin{equation}
u'_j=\widetilde u_j+0.45\,r_j b_j\tanh(v_j),\qquad 0\leq b_j\leq s.
\end{equation}
Free room excludes the prescribed minimum gap, preventing crossing under
simultaneous updates. Unlike the historical decoder, the anchored path has
no compulsory attraction toward uniformly spaced retained ranks.
Residual scale zero preserves proposal geometry.

\subsection{Final spline and deployment scope}
Selected coordinates become the interior of a clamped knot vector. One
standard endpoint-constrained B-spline least-squares refit produces the control
points. Deployment consists of one proposal/selection/decoder forward pass
and one final refit. Teacher search, BIC sweeps, Hard-Concrete sampling,
and greedy post-pruning are absent. Failed one-shot fits remain failures;
no hidden repair is used to make their reported results feasible.

\section{Training}
\subsection{Proposal and joint stages}
Proposal training emphasizes dense fit, parameter labels where available,
and ordered knot coverage. Joint training activates selection and
subset-conditioned fitting. Its initial geometry-calibration interval freezes
the encoder, ParameterHead, and proposal geometry while selection and the
subset decoder learn. Proposal geometry is subsequently unfrozen with a
reduced learning rate. Epoch counts are experiment settings, not fixed
properties of the method.

\subsection{Online counterfactual feasible teacher}
Training evaluates stochastic subsets, ranked-prefix alternatives, and bounded
counterfactual edits. A limited greedy search on selected training curves
seeks smaller feasible subsets. This is online numerical guidance, not an
offline cache, a separate pretrained teacher network, or a globally optimal
oracle.

Fixed-geometry deletion alone is insufficient: a changed mask also changes
the neural decoder. Teacher masks must therefore be decoded again and
fitted before acceptance as feasible mask/count labels. Geometric targets
include both parameters and knot positions. Useful but not yet redecoded-feasible
targets may supervise geometry; they do not become safe compact mask labels.
Teacher evaluations are detached targets. Ordinary derivatives through the
discrete deployed Top-$K$ are not claimed.

\subsection{Loss groups}
The objective combines fit/violation, local-fit and parameter terms, policy
and mask distillation, count calibration, boundary ranking, and geometric
distillation. Supervised source terms are masked out when labels are absent.
Complexity is encouraged within the configured feasibility-aware protocol,
not interpreted as proof of Equation~\ref{eq:target}.
\todo{Export final weights, search budgets, active objective configuration,
and checkpoint selection into an appendix. Do not replace these with an
inaccurate generic four-term loss.}

\subsection{Optional peak-error and depth studies}
For a separately declared squared-error target $\delta$, an optional loss
combines exact per-curve peak error and the mean of the largest $q$ fraction:
\begin{equation}
\mathcal L_{\rm peak}=\mathcal A_b\left[
\tfrac12\phi(\max_i e_{bi},\delta)+
\tfrac12\phi(\operatorname{meanTop}_{q}(e_{b:}),\delta)\right].
\end{equation}
$\phi$ is the configured normalized log-space fit penalty and $\mathcal A_b$
its batch tail-aware aggregation. The MSE pass criterion is unchanged.
Residual gradients do not make hard selection differentiable or guarantee
between-sample accuracy.

Optional proposal, selection, and survivor refinement depths all default to
zero extra layers. New residual gates start at zero, preserving the starting
function under compatible weight transfer, without guaranteeing improved
learning. Depth and peak-loss extensions have no complete measured results
in this draft.

\section{Experimental protocol}
\subsection{Data and supervision}
Synthetic training includes labelled splines and procedural shape families.
The audited runs use source $K=4,\ldots,24$, $M=192$, 1500 training curves
per epoch, and 500 synthetic validation curves. The new two-model study
trains general-purpose M16 and M32 models, with $K_c=16$ and $K_c=32$,
respectively. Both use mixed synthetic shapes, not dataset-specific models.
Their default labelled source ranges are $K=4,\ldots,16$ and
$K=4,\ldots,24$; this capacity-matched training changes the source
distribution as well as capacity. A capacity-only ablation must instead
use the shared source range $K=4,\ldots,16$ for both models.
The source minimality audit
is relative to tested source-subset removals, threshold, and parameterization;
it does not prove globally minimal free knots for finite noisy data.

External sources are UJI Pen Characters Version 2~\cite{uji2008},
Natural Earth coastlines~\cite{naturalearth}, and USGS
contours~\cite{usgscontours}. IndustrialOffset consists of procedurally
generated CAD-style offsets, not measured industrial observations.
Audited profiles use synthetic/procedural training and real validation;
writer or geographic groups separate validation from held-out tests.
The M16/M32 profiles retain this boundary: UJI, Natural Earth, and USGS
provide external validation, while IndustrialOffset remains test-only.
Both models are tested separately on all five sources, including the
same synthetic $K=4,\ldots,24$ evaluation range. Cases exceeding M16's
capacity are retained and identified rather than silently excluded.
External curves have no ground-truth spline knot vectors.

Input-fitting errors and original-reference-point errors are separate.
Upsampling short strokes does not create independent measured observations.
Normalization, resampling, groups, manifests, seeds, and content hashes
must accompany experiments.

\subsection{Six-method comparison}
Ours is compared with five disclosed adaptations: Park--Lee, Liang, Dung,
Kang, and Luo. All receive identical normalized points, the declared common
interior capacity, and an endpoint-constrained CPU float64 final refit.
Baselines use chord parameters while Ours predicts parameter residuals;
a fixed-chord ablation is needed to isolate knot selection.

Dung, Kang, and Luo may use a disclosed threshold-safe wrapper when their
native adapted result fails the common criterion. Bounded repairs and extra
refits remain in their results and latency. Native and wrapped settings must
be distinguished, and neither is an exact author implementation.

\subsection{Metrics, timing, and ablations}
Report pass fraction over all attempted curves, mean/P95 MSE, average
interior $K$, per-curve $\Peak$ mean/P95/maximum, and complete latency.
Exceptions and nonfinite outputs stay in the failure denominator.
Cardinality comparisons on jointly passing curves must also state sample
counts. Complete time includes parameterization, method execution, and
refit; network time is supplementary. GPU Ours versus CPU baselines is not
a hardware-neutral algorithmic speedup.

Controlled studies should include:
\begin{enumerate}
\item proposal/selection/survivor depth at fixed capacity, seed, data, and budget;
\item general-purpose M16 versus M32, with shared synthetic source ranges
for capacity-only ablations and separately disclosed capacity-matched training;
\item peak-loss on/off at the same MSE tolerance;
\item teacher geometry, redecoding verification, and parameter coupling.
\end{enumerate}
The default two-model experiment uses the same depth and peak-loss variant
for M16 and M32: two added refinement blocks per module and a peak-loss
weight of $0.05$ with squared-error target $5\times10^{-4}$.
The MSE tolerance remains $5\times10^{-5}$. Each model is initialized from
the same shallow K32 checkpoint; resizing candidate queries for M16
requires retraining and is not a function-preserving capacity reduction.
Within each six-method comparison all methods share that model's capacity;
M16 and M32 tables remain separate. No dataset router, test-error-based
model selection, or automatic M16-to-M32 fallback is introduced.
Explicit real-training-split adaptation would be a different protocol and
must be labelled separately. Test labels and test-fit outcomes cannot choose
capacity or model.

\section{Results and analysis --- to be completed}
\todo{Insert frozen paired six-method results, including synthetic strata and
all four external sources. No figures are required for this framework.}
\todo{Fill depth/capacity/peak-loss tables only after running the experiments;
no new improvement numbers currently exist.}

The 21 September audit is a cautionary diagnostic, not a final paper result.
On 122 identical inputs, Stable K32 and the saved Anchored 12+48 best model
achieved 60.66\% and 69.67\% pass rates, respectively, while mean knot counts
increased from 24.04 to 26.19. The latter was saved at epoch13, not epoch60.
This motivates assessing fit and compactness together. The repository audit
records conditional teacher denominators; these results must not be silently
merged with future formal experiments.

\section{Limitations and validity}
One-shot fits can fail the tolerance and retain too many knots. Candidate
capacity and proposal quality limit performance before selection. Teacher
verification covers only selected training curves; feasible targets do not
guarantee a feasible learned policy. Distribution shift, parameter bias,
and between-sample oscillation remain unresolved. No minimality theorem or
continuous error certificate is claimed.

Baseline adaptations show some severe original-reference instability despite
small input-grid residuals. Such failures need investigation and disclosure,
not silent removal. Procedural offsets do not establish general performance
on measured industrial geometry.

\section{Conclusion}
The framework combines candidate generation, adaptive one-shot selection,
coupled survivor geometry, and explicit B-spline fitting. Verified online
numerical guidance is used in training, without iterative subset search at
deployment. The final conclusion must follow the measured trade-off among
feasibility, complexity, peak error, and runtime, including unseen-source
behavior and failures.

\section*{Reproducibility and declarations}
\todo{Supply code/model hashes, manifests, seeds, commands, hardware, software,
training/inference budgets, and licenses.}
\todo{Supply authorship, funding, conflicts, and data availability declarations.}
\bibliographystyle{plainnat}
\bibliography{references}
\end{document}
